\section{Introduction}

Legged robots earn their mobility in exactly the settings that trip them up. A quadruped that can step over a curb or cross rubble is also one that can hook a foot under a loose cable, snag a calf on a taut wire, or drag a trailing line between its legs. We call this failure \emph{leg entanglement}, and on a walking platform such as the Unitree Go2~\cite{unitree_go2} it is both common and dangerous: the affected leg loses its expected swing, the controller keeps commanding a gait the limb can no longer execute, and the robot lurches, stalls, or falls. Detecting the condition early, and knowing which leg is caught, is a prerequisite for any sensible response.

Exteroception is a weak sensing route here: cables and wires are thin, dark, and usually occluded by the robot's own body, and the interaction happens below the belly, outside a forward camera's useful view. Proprioception is the natural alternative, since the joint encoders, torques, foot forces, and inertial signals that already drive the controller carry a direct mechanical trace of a snag~\cite{hwangbo2019learning,lee2020learning,camurri2017contact}. The trace, however, is subtle and easy to confuse. When a leg is held back the thigh joint fights a load it cannot move, so torque climbs while velocity stays near zero; but high thigh torque at low velocity also describes a robot standing still under its own weight. The confusion is sharpest for the rigid ``Lock'' stance the Go2 holds right after standing up, where torques are large and motion is nil yet nothing is wrong. A useful detector therefore cannot threshold effort in isolation; it must read effort in the temporal context of the gait.

We frame the task as time-series event detection over the proprioceptive stream and ask three questions at every timestep: \emph{whether} any leg is entangled, \emph{which} of the four legs (front-right, front-left, rear-right, rear-left) it is, and \emph{how severe} the entanglement is. Figure~\ref{fig:system} shows how a single causal model answers all three from a short window of on-robot signals, published as a state message on ROS~2~\cite{macenski2022ros2}. Answering the three jointly lets a shared representation carry the temporal context that separates entanglement from standing and from the Lock stance.

This paper makes the following contributions:
\begin{itemize}
  \item A proprioception-only formulation of leg entanglement as a multi-task, per-timestep event-detection problem that jointly answers \emph{whether}, \emph{which}, and \emph{how-severe}, using only signals the walking controller already produces.
  \item A compact causal dilated TCN with three heads and physics-grounded features that encode the ``high torque, little motion'' snag signature and separate entanglement from standing and from the Lock stance.
  \item A physics-grounded, calibrated per-leg severity index, reported as a severity index rather than a measured force, and an evaluation under a leakage-safe, whole-recording split with leave-one-recording-out cross-validation, a matched-protocol rule-based baseline, and a real-time ONNX runtime whose output matches the research pipeline.
\end{itemize}
